% ---
% Reemplaza líneas 1721-1760 de Docs/Technical_Report.tex
% Decisiones: D-033, H-038
% ---

\subsubsection{\textit{Kernel-Target Alignment} (KTA)}

El \textit{Kernel-Target Alignment} (KTA) cuantifica la similitud entre una matriz de
kernel $K \in \mathbb{R}^{N \times N}$ evaluada sobre $N$ muestras y la matriz de etiquetas
ideales de clasificación $Y = \mathbf{y}\mathbf{y}^\top$, donde $\mathbf{y} \in \{-1, +1\}^N$
denota el vector de etiquetas binarias de patología.

La formulación original propuesta por Cristianini et al. \cite{cristianini2002kernel}
define la alineación no centrada como el coseno de Frobenius entre ambas matrices:
\begin{equation}
  \widehat{\text{KTA}}_u(K, \mathbf{y}) =
  \frac{\langle K, \mathbf{y}\mathbf{y}^\top \rangle_F}{\|K\|_F \cdot \|\mathbf{y}\mathbf{y}^\top\|_F} =
  \frac{\mathbf{y}^\top K \mathbf{y}}{N\,\|K\|_F},
  \label{eq:kta_uncentered}
\end{equation}
donde $\langle A, B \rangle_F = \operatorname{Tr}(A^\top B) = \sum_{i,j} A_{ij} B_{ij}$ es el
producto interno de Frobenius y $\|\cdot\|_F$ su norma inducida. Dado que $y_i \in \{-1, +1\}$,
la norma de la matriz de etiquetas ideales satisface $\|\mathbf{y}\mathbf{y}^\top\|_F = N$.

No obstante, la métrica no centrada presenta un sesgo ante conjuntos de datos con
desbalance de clases. Si la proporción de clases difiere del 50\,\%, el valor esperado de
$y_i y_j$ es positivo, lo que provoca que una matriz de kernel no informativa
(e.g., una matriz constante u homogénea) obtenga valores elevados
de alineación simplemente por capturar la media de las etiquetas. Para suprimir este artefacto,
Cortes, Mohri y Rostamizadeh \cite{Cortes2012} introdujeron la \textit{alineación de kernel centrada}
(\textit{Centered Kernel Alignment}), adoptada en este trabajo como la métrica principal de KTA:
\begin{equation}
  \widehat{\text{KTA}}_c(K, \mathbf{y}) =
  \frac{\langle K_c, (\mathbf{y}\mathbf{y}^\top)_c \rangle_F}{\|K_c\|_F \cdot \|(\mathbf{y}\mathbf{y}^\top)_c\|_F} =
  \frac{\mathbf{y}^\top K_c \mathbf{y}}{\|K_c\|_F\,\sum_{i=1}^N (y_i - \bar{y})^2},
  \label{eq:kta_centered}
\end{equation}
donde $K_c = H K H$ representa la matriz de kernel centrada mediante la matriz de proyección
$H = I - \frac{1}{N}\mathbf{1}\mathbf{1}^\top$, con componentes escalares:
\begin{equation}
  (K_c)_{ij} = K_{ij} - \frac{1}{N}\sum_{k=1}^N K_{ik} - \frac{1}{N}\sum_{l=1}^N K_{lj} + \frac{1}{N^2}\sum_{k,l=1}^N K_{kl}.
  \label{eq:kernel_centering_scalar}
\end{equation}
El término de normalización $\|(\mathbf{y}\mathbf{y}^\top)_c\|_F$ coincide con la
suma de desviaciones cuadráticas de las etiquetas centradas, $\sum_{i=1}^N (y_i - \bar{y})^2$.

A pesar del centrado, los kernels que presentan concentración exhiben una cota
inferior analítica fija. Para la matriz identidad $K = I$ (el kernel de un espacio donde todas
las muestras son mutuamente ortogonales), la KTA centrada adopta un valor
constante e independiente de las etiquetas:
\begin{equation}
  \widehat{\text{KTA}}_c(I, \mathbf{y}) = \frac{1}{\sqrt{N - 1}},
  \label{eq:kta_identity_floor}
\end{equation}
que para la submuestra de evaluación de $N = 200$ lesiones equivale a $1/\sqrt{199} \approx 0.0709$.
(Para la formulación no centrada, dicho valor equivale a $1/\sqrt{N} = 1/\sqrt{200} \approx 0.0707$).
En consecuencia, las magnitudes absolutas de KTA no son directamente comparables entre kernels
cuyos espectros de autovalores difieren mucho. Para establecer una comparación válida,
la alineación de cada kernel se evalúa calculando el \textit{exceso de alineación} respecto a
la media de una prueba de permutación con $1\,000$ barajadas de las etiquetas que
conservan el número de lesiones de cada clase:
\begin{equation}
  \Delta\text{KTA}_c = \widehat{\text{KTA}}_c(K, \mathbf{y}) - \mathbb{E}_{\pi}\left[\widehat{\text{KTA}}_c(K, \mathbf{y}_\pi)\right].
  \label{eq:kta_excess}
\end{equation}
Para C4 y C5 el kernel es el de fidelidad (ecuación~\eqref{eq:quantum_kernel}); la referencia clásica es el RBF con $\gamma$ fijado por la heurística de la mediana; sobre los vectores de cada condición se calcula además la KTA del kernel lineal.
